\documentclass[11pt]{article}
\usepackage[margin=0.86in]{geometry}
\usepackage{amsmath,amssymb,mathtools,booktabs,microtype,xcolor,fancyhdr,hyperref}
\definecolor{navy}{HTML}{17365D}\definecolor{teal}{HTML}{147D78}
\hypersetup{colorlinks=true,linkcolor=navy,urlcolor=teal}
\pagestyle{fancy}\setlength{\headheight}{14pt}\fancyhf{}
\lhead{Scalar cyclic braid images}\rhead{Proof and verification}\cfoot{\thepage}
\newcommand{\Sgm}{\mathcal S}
\title{\color{navy}\bfseries Scalar Cyclic Finite Braid Images\\
\large Exact sigma--MGF coefficients in dimensions one through four}
\author{Computational verification note}\date{17 July 2026}
\begin{document}\maketitle

\begin{abstract}
We prove the finite-group sigma--MGF formula directly and test it on the
scalar cyclic quotients of $B_3$.  These examples isolate the scalar-center
selection rule while permitting arbitrary representation dimension.  Explicit
Reynolds-projector ranks agree with both the elementwise power-trace formula
and its conjugacy-class grouping in all tested dimensions $D=1,2,3,4$.
\end{abstract}

\section{The representation and the proposed quantity}
Fix $m\geq2$, $D\geq1$, and $\zeta=e^{2\pi i/m}$.  Define
\[
 \rho(\sigma_1)=\rho(\sigma_2)=\zeta I_D.
\]
The matrices are unitary and satisfy
$\rho(\sigma_1)\rho(\sigma_2)\rho(\sigma_1)
=\rho(\sigma_2)\rho(\sigma_1)\rho(\sigma_2)$, so this is a
$D$-dimensional representation of $B_3$.  Its image is $C_m$.

For a tuple $\tau=(\tau_1,\ldots,\tau_\ell)$ let
\[
 W_\tau=\bigotimes_{a=1}^{\ell}\operatorname{Sym}^{\tau_a}(\mathbb C^D),
 \qquad |\tau|=\sum_a\tau_a.
\]
The coefficient under test is
\[
 [m_\tau]\Sgm_{C_m}(\mathbf t),\qquad
 \Sgm_{C_m}(\mathbf t)=\frac1m\sum_{j=0}^{m-1}
 \prod_{a\geq1}\det(I-t_a\zeta^jI_D)^{-1}.
\]

\section{Proof of the formula and closed answer}
For any finite $G\leq U(V)$, the coefficient of $t^r$ in
$\det(I-tg)^{-1}$ is the character of $\operatorname{Sym}^rV$ at $g$.
Multiplying the factors and applying the Reynolds character formula gives
\[
 [m_\tau]\Sgm_G
 =\frac1{|G|}\sum_{g\in G}\prod_a
 \operatorname{Tr}(g\mid\operatorname{Sym}^{\tau_a}V)
 =\dim(W_\tau^G).
\]
Furthermore,
\[
 \det(I-tg)^{-1}=\exp\left(\sum_{k\geq1}
 \frac{t^k}{k}\operatorname{Tr}(g^k)\right).
\]
Grouping the last average into conjugacy classes yields the weight
$|[c]|/|G|=1/|C_G(c)|$, proving both proposed finite-image formulas.

In the present representation, $\zeta^jI_D$ acts on all of $W_\tau$ by
$\zeta^{j|\tau|}$.  Hence the Reynolds operator is either zero or identity:
\[
 \boxed{[m_\tau]\Sgm_{C_m}=\mathbf1_{m\mid|\tau|}
 \prod_a\binom{D+\tau_a-1}{\tau_a}.}
\]
This proves the scalar selection rule and also determines the nonzero value,
not merely its support.

\section{Independent verification strategy}
The program first closes the matrix group from the displayed generator and
asserts order $m$, unitarity, and the braid relation.  It then constructs an
orthonormal occupation basis of every $\operatorname{Sym}^{r}\mathbb C^D$
inside $(\mathbb C^D)^{\otimes r}$.  In that basis it forms
\[
 P_\tau=\frac1m\sum_{g\in C_m}\bigotimes_a\operatorname{Sym}^{\tau_a}(g)
\]
and computes its rank.  This is the direct target calculation.

The comparator never uses this projector.  It obtains the complete
homogeneous characters from Newton's recurrence
\[
 r h_r(g)=\sum_{k=1}^r\operatorname{Tr}(g^k)h_{r-k}(g),\qquad h_0=1,
\]
then averages $\prod_a h_{\tau_a}(g)$ elementwise and, separately, by
matrix conjugacy classes.

\begin{center}
\begin{tabular}{@{}rrrrrr@{}}\toprule
$m$&$D$&$\tau$&$\dim W_\tau$&$\operatorname{rank}P_\tau$&formula\\\midrule
3&1&$(1)$&1&0&0\\
3&1&$(3)$&1&1&1\\
3&2&$(2,1)$&6&6&6\\
4&2&$(4)$&5&5&5\\
4&3&$(2,2)$&36&36&36\\
5&4&$(5)$&56&56&56\\\bottomrule
\end{tabular}
\end{center}
The maximum projector idempotence/Hermiticity error in this sweep is below
$6\times10^{-15}$.

\section{Reproduction and scope}
From the repository root run
\begin{verbatim}
.venv/bin/python lambda_distributions/proofs/finite_braid_images/
  cyclic_scalar/verification.py
.venv/bin/marimo edit marimo_notebooks/braid_cyclic_scalar.py
\end{verbatim}
after joining each wrapped command.  The proof is valid for every $m,D,\tau$;
the table is a small numerical regression suite, not evidence substituted for
the proof.
\end{document}

